\chapter{Introduction}
\label{sec:introduction}

% \begin{flushright}
% \begin{minipage}{0.6\textwidth}
% \raggedleft
% \emph{''I learned very early the difference between knowing the name of something and knowing something.''}\\
% \textemdash{} Richard Feynman
% \end{minipage}
% \end{flushright}

Understanding what neural networks compute, not only what they output,
is the central problem of mechanistic interpretability. A model that correctly solves arbitrary problems or chains
reasoning steps coherently need not be implementing anything resembling the algorithm a human would use.
Whether it does, where in the network that computation lives, and whether it can be modified in controlled ways
are questions that evaluation benchmarks cannot answer.

\section{Mechanistic Interpretability}

As neural networks scale from idealised models to agentic systems executed in consequential and complex settings,
the gap between performance and understanding widens. A model achieving near-human accuracy on reasoning benchmarks
need not be implementing anything resembling the intended algorithm. A system that has learned approximate heuristics is indistinguishable from one that has learned the correct
procedure until it fails on inputs the heuristic does not cover.

Early interpretability work was largely descriptive, using saliency maps, attribution scores, and probing classifiers. In vision models, this approach identified low-level features such as edge and curve detectors that recur across architectures and seeds \citep{olah2020zoom}, while analysis of transformer attention revealed heads implementing specific operations, such as induction heads completing repeated sequences \citep{elhage2021mathematical}. The grokking phenomenon offers a clearer demonstration: small models trained on modular arithmetic first memorise their training data, then undergo a phase transition to a compressed, generalising algorithm \citep{power2022grokking}. Fourier analysis shows this algorithm is implemented via trigonometric identities in the residual stream, bearing no surface resemblance to the training objective \citep{nanda2023progress}, a result extended by further work on arithmetic transformers \citep{zhong2023clock, quirke2024understanding, quirkeneo2025understanding, musat2024clustering}.
These findings establish that small models converge on interpretable algorithms independent of surface statistics, but manual circuit analysis does not scale to large language models, which are trained on far more diverse data without explicit algorithmic curricula. Whether their structured reasoning reflects stable, localisable computation or superficial association remains difficult to establish.


Anthropic's \textit{On the Biology of a Large Language Model}~\cite{lindsey2025biology} proposes
a reproducible foundational methodology for approaching interpretability on an LLM.
Combining sparse autoencoders (SAEs) and backward-pass attribution,
it constructs dependency graphs that trace information flow from input tokens through intermediate
features to output logits. This approach forms the methodological foundation of this thesis, which
examines how structured reasoning is implemented within large instruction-tuned models and what
the geometry of those representations reveals about whether model behaviour is grounded in stable, localisable algorithms.

\section{Notation and Methods}
\label{sec:interp-methods}
Let $\mathcal{M}$ be a transformer with $L$ layers, residual stream dimension $d$, and MLP
sublayers $\mathrm{MLP}^{(0)}, \ldots, \mathrm{MLP}^{(L-1)}$.  All methods employed in this
thesis function on the residual stream of $\mathcal{M}$.  Let $\mathbf{h}^{(\ell)}_t \in \mathbb{R}^d$
denote the residual-stream vector at layer $\ell \in \{0, \ldots, L{-}1\}$ and token position $t$.
The contribution of each layer is additive:
\begin{equation}
    \mathbf{h}^{(\ell+1)}_t = \mathbf{h}^{(\ell)}_t
      + \mathrm{Attn}^{(\ell)}\!\left(\mathbf{H}^{(\ell)}\right)_t
      + \mathrm{MLP}^{(\ell)}\!\left(\mathbf{h}^{(\ell)}_t + \mathrm{Attn}^{(\ell)}\!\left(\mathbf{H}^{(\ell)}\right)_t\right),
    \label{eq:residual}
\end{equation}
where $\mathbf{H}^{(\ell)} \in \mathbb{R}^{T \times d}$ denotes the entire sequence of residual-stream
vectors at layer $\ell$, with $T$ the sequence length.
This additive structure makes the residual stream appropriate for linear analysis, as each
layer writes into a shared representation that accumulates with layers.

Mechanistic analysis replaces each MLP sublayer $\mathrm{MLP}^{(\ell)}$ with a pretrained
\textit{transcoder} $\mathcal{T}_\ell$, a sparse encoder--decoder layer whose weights are trained separately 
to reconstruct the MLP output while expressing it as a sparse combination of \(K\) learned feature directions~\citep{bricken2023monosemanticity, cunningham2023sparse}.  
For the input $\mathbf{x} \in \mathbb{R}^d$ and where $\mathbf{W}_\mathrm{enc}^{(\ell)} \in \mathbb{R}^{d \times K}$  to the MLP sublayer at layer $\ell$, the transcoder computes
the following pre-activations
\begin{equation}
  z_{\ell,k}(\mathbf{x})
    = \bigl(\mathbf{W}_\mathrm{enc}^{(\ell)}\,\mathbf{x}
           + \mathbf{b}_\mathrm{enc}^{(\ell)}\bigr)_k,
  \label{eq:preact}
\end{equation}
so that feature $k$ has activation
$a_{\ell,k}(\mathbf{x}) = \mathrm{ReLU}(z_{\ell,k}(\mathbf{x}))$,
and the reconstructed MLP output is
\begin{equation}
  \hat{f}_\ell(\mathbf{x})
    = \mathbf{W}_\mathrm{dec}^{(\ell)}\,\mathbf{a}_\ell(\mathbf{x})
    + \mathbf{b}_\mathrm{dec}^{(\ell)},
  \label{eq:recon}
\end{equation}
where $\mathbf{W}_\mathrm{dec}^{(\ell)} \in \mathbb{R}^{K \times d}$ and
$\mathbf{a}_\ell(\mathbf{x})$ is the vector of all feature activations at layer $\ell$.
Feature sparsity implies that at most a
small active set of features contributes for any given input, allowing for more interpretable analysis and monosemanticity.

\vspace{3 mm}
\textit{Attribution graphs}. Attribution scores are computed on the \textit{linearised} model, in which attention outputs
are disconnected and treated as constants, so that residual stream propagates only through the
transcoder reconstruction path.  Due to linearisation, the attribution scores
are additive and interpretable as influence flows through the network.  Let \(T\) denote the position of the final input token, and let
\(\mathbf{h}_T \in \mathbb{R}^d\) be the corresponding pre-unembedding hidden
state. Let \(V\) denote the model vocabulary and
\(W_U \in \mathbb{R}^{d \times |V|}\) the unembedding matrix. For output token
\(j \in V\), the unembedding direction is
\begin{equation}
    \mathbf{v}_j
    =
    W_U[:,j]
    -
    \frac{1}{|V|}\sum_{r \in V} W_U[:,r].
    \label{eq:demeaned-unembedding}
\end{equation}
The target logit is then
\begin{equation}
    y_j = \mathbf{v}_j^\top \mathbf{h}_T ,
    \label{eq:target-logit}
\end{equation}


and the attribution score
of feature $k$ at position $t$ in layer $\ell$ is
\begin{equation}
    \alpha_{\ell,k,t}
      = a_{\ell,k,t}
        \cdot \frac{\partial y_j}{\partial a_{\ell,k,t}},
    \label{eq:attribution}
\end{equation}
the product of the feature activation with the gradient of the target logit with respect to
that activation~\citep{lindsey2025biology}.  Obtaining all such scores over a given prompt
yields a directed weighted graph, the \textit{attribution graph}, in which nodes are
(layer, position, feature) triples and edges hold the influence $\alpha_{\ell,k,t}$ from
earlier features to later ones.  The graph is reduced in size by retaining only those nodes and edges that account for a fixed fraction of the total downstream influence on $y_j$, given a
a node threshold $\tau_\mathrm{node}$ and an edge threshold $\tau_\mathrm{edge}$.

\textit{Activation patching}. For verifying causal hypotheses~\citep{meng2022locating, conmy2023automated}, given an original prompt
$x$ and a counterfactual prompt $x'$, patching replaces the residual stream at layer $\ell$
and position $t$,
\begin{equation}
    \tilde{\mathbf{h}}^{(\ell)}_t = \mathbf{h}'^{(\ell)}_t,
    \label{eq:patching}
\end{equation}
and measures the resulting change in output logit,
$\Delta_j = y_j(\tilde{\mathbf{h}}) - y_j(\mathbf{h})$.  A non-zero $\Delta_j$ shows
that the patched representation is sufficient to alter the target logit. 

\textit{Steering vectors} are extracted by contrasting two prompt sets,
following representation-engineering and contrastive activation-addition
methods~\citep{zou2023representation, rimsky2023steering}.
Let \(D_+\) denote prompts where a concept is present and \(D_-\) prompts
where it is absent. The steering direction at layer \(\ell\) is then the
mean residual-stream difference at a designated anchor position \(t^*\):
\begin{equation}
    \delta^{(\ell)} =
    \mathbb{E}_{x \in D_+}\!\left[\mathbf{h}^{(\ell)}_{t^*(x)}\right]
    -
    \mathbb{E}_{x \in D_-}\!\left[\mathbf{h}^{(\ell)}_{t^*(x)}\right].
    \label{eq:delta}
\end{equation}

Under the linear representation hypothesis~\citep{park2023linear}, $\delta^{(\ell)}$
approximates the axis along which the concept is encoded at depth $\ell$. This
one-dimensional summary is a useful first-order model, although not all language-model
features are expected to be strictly one-dimensional~\citep{engels2024notall}. At inference
time, the model can be steered by adding $\alpha_{\mathrm{steer}}\,\delta^{(\ell)}$ to
the residual stream, where $\alpha_{\mathrm{steer}} \in \mathbb{R}$ controls the strength
and direction of the intervention. This shifts the model's output without modifying any weights.

\section{Structure of Thesis}

This thesis studies whether mechanisms described in Anthropic's \textit{On the Biology of a Large Language Model} can be identified in the open-access model \texttt{Qwen3-4B-Instruct}. Its first aim is to reproduce two of the behaviours analysed in that work under an independent implementation: the arithmetic addition circuit, with particular attention to carry-related lookup features and attribution graphs, and a multilingual circuit experiment testing whether operation, operand, and output-language components can be separately intervened on. These reproductions are treated as mechanistic tests of whether analogous sparse features, attribution-graph structures, and causal intervention effects transfer to a different model family.

The second aim is to extend this analysis through a systematic concept-localisation study. These experiments ask how scientific and numerical concepts emerge as the model processes a prompt: whether the relevant direction appears gradually or as a sharp layer-wise spike, whether it stabilises across depth, and how narrowly it can be localised to particular layers, token positions, or sparse features. Using contrastive residual-stream directions, feature alignment, activation grids, and targeted interventions, the thesis compares concepts across three modular arithmetic problems: carry propagation in integer addition, GCD divisibility, and residue class membership.

The third aim is to examine whether the internal representations identified through concept localisation can be exploited for targeted knowledge editing. Using soft prompting and prefix tuning, this work tests whether a small set of learned continuous vectors, prepended to the input while keeping all base model parameters frozen, is sufficient to shift model behaviour on concepts that prove difficult to localise cleanly. The approach is evaluated across three tasks of varying difficulty: residue class membership, balanced parentheses, and causal direction judgements, with results compared against the base model to assess the extent to which soft interventions can compensate for the absence of a sharp, geometrically localised concept representation.

The rest of this dissertation is structured as follows. Chapter~\ref{sec:related} provides the relevant background, covering mechanistic interpretability, arithmetic circuits in transformers, and concept localisation methods. Chapter~\ref{sec:repro} presents a reproducibility study of the addition circuit and multilingual experiment from Lindsey et al., examining whether the same mechanistic signatures appear in \texttt{Qwen3-4B-Instruct}. Chapter~\ref{sec:concept_localisation} extends this analysis through a systematic concept-localisation study, tracing how arithmetic concepts emerge across layers, token positions, and sparse features using contrastive residual-stream directions. Chapter~\ref{sec:k_e} examines knowledge editing methodologies, exploring soft prompting and prefix tuning as targeted interventions on model representations. Finally, Chapter~\ref{sec:conclusion} concludes the dissertation and outlines directions for future work.